\begin{definition}[The per-process variables beside the ABA network]
  \label{def:aba-state}
  \lean{PLTS.ABA.ABAState}\lean{PLTS.ABA.ABAState.processes}\lean{PLTS.ABA.ABAState.decidedSent}\lean{PLTS.ABA.ABAState.decidedReceived}\lean{PLTS.ABA.ABAState.corrupted}\lean{PLTS.ABA.ABAState.F}\lean{PLTS.ABA.ABAState.startedRounds}\lean{PLTS.ABA.ABAState.decidedCount}\lean{PLTS.ABA.ABAState.setProcessVariables}\lean{PLTS.ABA.ABAState.stepRound}\lean{PLTS.ABA.ABAState.sendDecidedOnGrade2}\lean{PLTS.ABA.ABAState.recordDecided}\lean{PLTS.ABA.ABAState.deliverDecided}\lean{PLTS.ABA.ABAState.corrupt}\leanok
  \uses{def:aba-labels, def:aba-round-loop-state, def:aba-programs, def:aba-composed}
  The round loop is carried by two components, and $\mathrm{ABAState}$ reads them as one object:
  \[
    \mathrm{ABAState}(P) \;=\;
      \bigl([n] \to \mathrm{RoundLoopVariables}(n)\bigr) \times \mathrm{ABANetworkState}(n),
  \]
  the round-loop variables of each process --- the control state of Definition~\ref{def:aba-round-loop-state} together with its DECIDED received sets and its own record of the DECIDED bits it has multicast --- beside the ABA network of Definition~\ref{def:aba-composed}, which holds the DECIDED sets, the corrupted set with its budget and the rounds a correct process has started (D37).
  Six accessors gather what the two components hold apart.
  $\mathrm{processes}$ reads a process's control state, $\mathrm{decidedReceived}$ its received sets and $\mathrm{corrupted}$ the flag saying whether its program has been replaced (D23), all three at that process's variables; $\mathrm{decidedSent}$ reads the multicast sent sets, $F$ the corrupted set and $\mathrm{startedRounds}$ the started rounds, all three at the network.
  A process's own record of its DECIDED sends is the field $decidedSent$ of its round-loop variables, which its relay and its return read, and which agrees with the network's at every process outside $F$ (D12$'$).
  The flag and the corrupted set live in different components, and the two are written by the two halves of one $\mathrm{fail}$ transition; that they name the same processes is the first conjunct of the core simulation's invariant and not a fact of the state shape.
  Read through them, $\mathrm{decidedCount}\,i\,b$ is the number of distinct senders whose $\langle\mathrm{DECIDED},b\rangle$ has been delivered to $i$, which is the count the return guard reads, and the operations $\mathrm{setProcessVariables}$, $\mathrm{stepRound}$, $\mathrm{sendDecidedOnGrade2}$, $\mathrm{recordDecided}$, $\mathrm{deliverDecided}$ and $\mathrm{corrupt}$ name the effects of the transitions below.
  The coin return's effect is $\mathrm{stepRound}$ and the DECIDED send's is $\mathrm{sendDecidedOnGrade2}$.
  An accessor is a projection of one component, so nothing of the two components is identified: the object is the pair read through accessors and not a coarser state.
  Full construction: \leandecl{PLTS.ABA.ABAState}.
\end{definition}
